% Generated draft: requires mathematical and source review.
\documentclass[11pt]{article}
\usepackage[margin=25mm]{geometry}
\usepackage{fontspec}
\setmainfont{Latin Modern Roman}
\usepackage{amsmath,amssymb,mathtools,unicode-math}
\setmathfont{Latin Modern Math}
\usepackage{enumitem,needspace,xurl,hyperref}
\setlength{\parindent}{0pt}
\setlength{\parskip}{5pt}
\setlength{\emergencystretch}{3em}
\newcommand{\N}{\mathbb{N}}
\newcommand{\Z}{\mathbb{Z}}
\newcommand{\Q}{\mathbb{Q}}
\newcommand{\R}{\mathbb{R}}
\newcommand{\C}{\mathbb{C}}
\newcommand{\F}{\mathbb{F}}
\newcommand{\A}{\mathbb{A}}
\newcommand{\PP}{\mathbb{P}}
\newcommand{\E}{\mathbb{E}}
\newcommand{\eps}{\varepsilon}
\newcommand{\dd}{\,\mathrm{d}}
\begin{document}
\section*{UnsolvedMath ID: 8700060}
\textbf{Source problem code:} catalogue entry
\subsection{Definitions and mathematical statement}
Let $A/\mathbb Q$ be a simple abelian variety of dimension $g\geq1$. Let $A(\mathbb R)^0$ be the identity component of its real points and let
\[
\exp_A:\mathbb R^g\longrightarrow A(\mathbb R)^0,\qquad
\ker\exp_A=\Omega=\mathbb Z\omega_1+\cdots+\mathbb Z\omega_g.
\]
For $u=\sum_{i=1}^g u_i\omega_i$ with $u_i\in\mathbb R$, assume $\exp_A(u)\in A(\mathbb Q)$. The source asserts that $1,u_1,\ldots,u_g$ are linearly independent over $\mathbb Q$, that is,
\[
a_0+\sum_{i=1}^g a_i u_i=0,\quad a_i\in\mathbb Q
\quad\Longrightarrow\quad a_0=\cdots=a_g=0.
\]
The coordinates are relative to the displayed integral period-lattice basis.

As written, the assertion has a missing necessary qualification: $u=0$ satisfies the premise but gives zero coordinates, so the asserted independence fails. More generally, if $\exp_A(u)$ is torsion then some positive integer $q$ has $qu\in\Omega$, making every $u_i$ rational. Thus a non-torsion restriction is necessary for the intended independence claim, but it is not present in the supplied statement and is not silently inserted. The exact corrected source conjecture requires verification.
\subsection{Short English statement}
The source predicts rational linear independence of one and the period-basis coordinates of a real abelian logarithm of a rational point. Its statement omits a necessary non-torsion condition, so the literal zero-logarithm case already prevents the claimed conclusion.
\subsection{Sources}
\begin{itemize}
\item Michel Waldschmidt, \emph{Open Diophantine Problems} (2004), Conjecture 4.18, \url{https://arxiv.org/abs/math/0312440}; accessed 2026-07-29, source-TeX extraction.
\item The supplied research summary mentions Bertrand's work and abelian/elliptic logarithm independence as related context without complete citations or precise theorem hypotheses.
\item Local registry \texttt{data/unsolvedmath/problems.json}, permanent UnsolvedMath ID $8700060$, source code catalogue entry, \emph{Conjecture 4.18}. Dataset attribution: Ulam AI and the UnsolvedMath Contributors.
\end{itemize}
\end{document}
